
Multi-objective code generation architectures assume quality dimensions can be factorized into separable subspaces. This work presents a methodological framework to test this assumption through spectral analysis of code modification outcomes. We validate the framework on synthetic test data (10,000 commits), demonstrating that the methodology correctly identifies when factorization structure is absent. The synthetic data exhibited statistical independence (cross-aspect coupling = 0.072 ≤ 0.2) but no aspect-dominant directional structure (spectral gap λ₁/λ₄ = 1.580 < 2.0, permutation test p = 0.955). This proof-of-concept establishes that independence does not imply factorization—covariance geometry can be spherical rather than aspect-aligned. The validation framework combines residual covariance analysis, spectral decomposition, permutation testing (1000 iterations), directional stability measures, and cross-validation. We acknowledge the critical limitation: all results use synthetic test data for pipeline validation, not real GitHub commits. Real data collection (Phase 1A: 5 days commit mining, Phase 1B: 22 hours metric computation) is required before scientific conclusions about developer behavior can be drawn. The methodological contribution provides researchers with validation tools to test factorization assumptions empirically before architectural commitment.
